\chapter{Análisis económico y de mercado}
\label{cap:economico}

El proyecto fue financiado íntegramente por el autor, sin auspiciantes ni
aportes de la Universidad. Este capítulo resume lo que costó desarrollarlo y
dónde podría ubicarse el resultado. El análisis es deliberadamente breve: la
cifra más significativa no es un monto, sino que el costo de herramientas fue
nulo.

\section{Costo de materiales}

El hardware se reduce a la placa de desarrollo y a una caja para protegerla
durante el trabajo (tabla~\ref{tab:costo-materiales}). No se fabricó ninguna
placa propia.

\begin{table}[ht]
  \centering
  \begin{tabular}{lr}
    \toprule
    \textbf{Ítem} & \textbf{Costo (USD)} \\
    \midrule
    PolarFire SoC Discovery Kit & 128 \\
    Caja estanca de protección & 5 \\
    \midrule
    \textbf{Total} & \textbf{133} \\
    \bottomrule
  \end{tabular}
  \caption{Costo de materiales del prototipo.}
  \label{tab:costo-materiales}
\end{table}

\section{Costo de herramientas y licencias}

Cero. Todas las herramientas libres del capítulo~\ref{cap:metodologia} son
gratuitas, y la única herramienta imprescindible sin alternativa libre,
Libero SoC, se usó con la licencia Silver, que Microchip otorga sin costo
para los dispositivos de menor capacidad (sección~\ref{sec:licencias}). Los
servicios de alojamiento del repositorio, de las versiones y de la
demostración pública también se usaron en sus modalidades gratuitas.

Este resultado no es un detalle: es lo que hace reproducible el proyecto.
Cualquier estudiante con una placa de USD 128 puede repetir el desarrollo
completo, desde la simulación hasta la liberación de una versión, sin pagar
una sola licencia. Un producto comercial que requiriera un dispositivo más
grande sí necesitaría una licencia paga de Libero.

\section{Esfuerzo de ingeniería}

El esfuerzo se estima en unas dos horas por día hábil durante un año, es
decir, del orden de \textbf{490 horas} de trabajo, a cargo de una sola
persona.

\section{Análisis de mercado}

\paragraph{El contexto}
MPEG-2 es un estándar maduro. Los códecs posteriores lo superan en
eficiencia de compresión, y cualquier procesador moderno lo decodifica por
software sin esfuerzo. Pero sigue presente en la televisión digital de
varios países, en el DVD, en equipos de difusión y en grandes archivos de
video. Y, desde el vencimiento de sus patentes en 2018
(sección~\ref{sec:patentes}), es el único códec de video de uso masivo que
puede implementarse en hardware sin pagar regalías en casi todo el mundo.

\paragraph{Dónde tiene sentido este producto}
No compite con los decodificadores integrados en circuitos de consumo, que
cuestan centavos a escala, ni con la decodificación por software. Su lugar
está en tres nichos:

\begin{itemize}
  \item \textbf{Formación}: una plataforma abierta, documentada y verificada
    para enseñar decodificación de video en hardware, integración de un núcleo
    de propiedad intelectual con un procesador y el ciclo completo de un
    producto. Es el uso más natural, y el que motivó el proyecto.
  \item \textbf{Sistemas embebidos con FPGA}: equipos que ya tienen una FPGA
    y necesitan decodificar MPEG-2 con consumo bajo y comportamiento
    determinista, sin agregar un procesador dedicado al video.
  \item \textbf{Base para otros desarrollos}: un diseño de referencia libre
    sobre el que portar el núcleo a otras FPGA o extenderlo, con un proceso
    de verificación y de liberación ya resuelto.
\end{itemize}

\paragraph{Modelo}
Las licencias elegidas definen el modelo: el producto es abierto, y su valor
está en el diseño verificado, la documentación y el proceso, no en el
secreto. La licencia recíproca del hardware propio
(sección~\ref{sec:licencia-propia}) asegura que ese valor siga disponible
para quien lo reutilice.
